% formal/sys-landscape/pentagon-rotation-formula.tex — input by formal/main.tex
%
% Private proof draft for the explicit systolic-ratio formula of
% `P_5 x_L R(theta) P_5`.
%
% Sources:
%   - Generated data: experiments/sys-landscape/pentagon-rotation-formula/theta-sweep.jsonl
%   - HKO source: papers/hko2024/counterexample.tex
%   - Research log: research/sys-landscape/design/pentagon-rotation-formula.md
%
% Structure:
%   - Theorem statement for the explicit formula
%   - Proof skeleton and intermediate lemmas
%   - Explicit statement of gaps and verification targets

\subsection{Pentagon Rotation Formula (Draft)}
\label{sec:pentagon-rotation-formula}

This section is a private proof draft for the explicit formula of
\[
  \mathrm{sys}\!\left(P_5 \times_L R(\theta) P_5\right).
\]
It is intended for agent and Jörn review, not yet thesis-facing exposition.

Let
\[
  v_k := \left(
    \cos\left(\frac{2\pi k}{5}\right),
    \sin\left(\frac{2\pi k}{5}\right)
  \right),
  \qquad 0 \le k \le 4,
\]
and let~$P_5$ be the regular pentagon with those vertices. Since rotating both
factors by the same angle is symplectic, the experiment convention used in
`experiments/sys-landscape/pentagon-rotation-formula/main.rs` is symplectically
equivalent to the standard one used here. The HKO special case uses the second
pentagon with vertices at angles $-\pi/2 + 2\pi k/5$, i.e. a rotation by
$-\pi/2$. Modulo the pentagon's $2\pi/5$ rotational symmetry this is the same
midpoint configuration as the present parameter value $\theta=\pi/10$, up to
the sign convention for rotating the second factor.

\begin{unverified}
\paragraph{Conjectural target.}
Let~$P_5 \subset \mathbb{R}^2$ be the regular pentagon with circumradius~$1$
in the angle convention used by the experiment code. For
\[
  K_\theta := P_5 \times_L R(\theta) P_5,
  \qquad 0 \le \theta \le \frac{\pi}{5},
\]
one has
\[
  \mathrm{sys}(K_\theta)
  =
  \frac{5 + 2\sqrt{5}}{10}
  \sec^2\!\bigl(\min(\theta, \tfrac{\pi}{5} - \theta)\bigr).
\]
Equivalently,
\[
  \mathrm{sys}(K_\theta)
  =
  \frac{5 + 2\sqrt{5}}{10}\sec^2(\theta)
  \quad \text{for } 0 \le \theta \le \frac{\pi}{10},
\]
and the second half of the fundamental domain follows by the symmetry
$\theta \mapsto \frac{\pi}{5} - \theta$.
\end{unverified}

\begin{unverified}
\begin{proposition}[The active 2-bounce branch has the claimed secant law]
\label{prop:pentagon-rotation-two-bounce}
Let~$K_\theta := P_5 \times_L R(\theta)P_5$ with
$0 \le \theta \le \frac{\pi}{10}$. Among the 2-bounce trajectories joining a
vertex of~$P_5$ to the opposite edge, the minimizing branch is attained at
\[
  \lambda(\theta)
  =
  \frac{1}{2}
  -
  \frac{\tan(\theta)}{2\tan(\pi/10)},
\]
and its capacity is
\[
  c_{2\text{-bounce}}(K_\theta)
  =
  \left(1+\cos\left(\frac{\pi}{5}\right)\right)^2 \sec(\theta).
\]
Consequently,
\[
  \mathrm{sys}(K_\theta)
  \le
  \frac{5 + 2\sqrt{5}}{10}\sec^2(\theta)
  \qquad
  \text{for } 0 \le \theta \le \frac{\pi}{10}.
\]
\end{proposition}

% [TODO: JÖRN - confirm the exact preferred theorem statement once the private
% proof is stable.]

\begin{proof}
We compute the active 2-bounce branch explicitly on the half-domain
$0 \le \theta \le \pi/10$.

\paragraph{2-bounce family.}
The combinatorial classification of 2-bounce trajectories that cannot be
translated into the interior depends only on the table pentagon and is the same
as in the HKO proof: they are exactly the chords joining a vertex to a point on
the opposite edge. By cyclic symmetry it is enough to study the family joining
$v_0$ to the edge $[v_2,v_3]$.

\[
  x_\lambda := \lambda v_2 + (1-\lambda) v_3,
  \qquad 0 \le \lambda \le 1.
\]
Since
\[
  v_2 = \left(
    -\cos\left(\frac{\pi}{5}\right),
    \sin\left(\frac{\pi}{5}\right)
  \right),
  \qquad
  v_3 = \left(
    -\cos\left(\frac{\pi}{5}\right),
    -\sin\left(\frac{\pi}{5}\right)
  \right),
\]
we have
\[
  x_\lambda
  =
  \left(
    -\cos\left(\frac{\pi}{5}\right),
    (2\lambda-1)\sin\left(\frac{\pi}{5}\right)
  \right),
\]
and therefore the chord vector is
\[
  d_\lambda := x_\lambda - v_0
  =
  \left(
    -\left(1+\cos\left(\frac{\pi}{5}\right)\right),
    (2\lambda-1)\sin\left(\frac{\pi}{5}\right)
  \right).
\]

For the rotated support pentagon, write
\[
  w_k(\theta) := R(\theta)v_k.
\]
On the half-domain $0 \le \theta \le \pi/10$, the vector $-d_\lambda$ stays in
the normal cone of $w_0(\theta)$, while $d_\lambda$ stays in the union of the
adjacent normal cones of $w_2(\theta)$ and $w_3(\theta)$. Therefore
\[
  h_{R(\theta)P_5}(-d_\lambda)
  =
  \langle -d_\lambda, w_0(\theta)\rangle,
\]
and
\[
  h_{R(\theta)P_5}(d_\lambda)
  =
  \max\Bigl(
    \langle d_\lambda, w_2(\theta)\rangle,
    \langle d_\lambda, w_3(\theta)\rangle
  \Bigr).
\]
The support-switch expression
\[
  \langle d_\lambda, w_2(\theta)-w_3(\theta)\rangle
\]
is affine and strictly increasing in~$\lambda$, because
\[
  \frac{d}{d\lambda}
  \langle d_\lambda, w_2(\theta)-w_3(\theta)\rangle
  =
  4\sin^2\left(\frac{\pi}{5}\right)\cos\theta
  >
  0
  \qquad
  \text{for } 0 \le \theta \le \frac{\pi}{10}.
\]
Hence there is a unique switch point, characterized by
\[
  \langle d_\lambda, w_2(\theta)-w_3(\theta)\rangle = 0.
\]
Since
\[
  w_2(\theta)-w_3(\theta)
  =
  R(\theta)(v_2-v_3)
  =
  2\sin\left(\frac{\pi}{5}\right)(-\sin\theta,\cos\theta),
\]
the support-switch equation becomes
\[
  \left(1+\cos\left(\frac{\pi}{5}\right)\right)\sin\theta
  +
  (2\lambda-1)\sin\left(\frac{\pi}{5}\right)\cos\theta
  =
  0.
\]
Using
\[
  \frac{1+\cos(\pi/5)}{\sin(\pi/5)}
  =
  \frac{2\cos^2(\pi/10)}{2\sin(\pi/10)\cos(\pi/10)}
  =
  \cot(\pi/10),
\]
this yields
\[
  \lambda(\theta)
  =
  \frac{1}{2}
  -
  \frac{\tan(\theta)}{2\tan(\pi/10)}.
\]
In particular, $\lambda(0)=1/2$ gives the opposite-edge midpoint and
$\lambda(\pi/10)=0$ gives the diagonal endpoint $v_3$.

Let
\[
  L_\theta(\lambda)
  :=
  h_{R(\theta)P_5}(-d_\lambda)
  +
  h_{R(\theta)P_5}(d_\lambda)
\]
be the 2-bounce $T^\circ$-length along this family. Since
$\langle d_\lambda,w_2(\theta)-w_3(\theta)\rangle$ is increasing in~$\lambda$,
the active support is~$w_3(\theta)$ for $\lambda \le \lambda(\theta)$ and
$w_2(\theta)$ for $\lambda \ge \lambda(\theta)$. Therefore
\[
  L_\theta(\lambda)
  =
  \langle -d_\lambda, w_0(\theta)\rangle
  +
  \langle d_\lambda, w_3(\theta)\rangle
  \qquad
  \text{for } \lambda \le \lambda(\theta),
\]
and
\[
  L_\theta(\lambda)
  =
  \langle -d_\lambda, w_0(\theta)\rangle
  +
  \langle d_\lambda, w_2(\theta)\rangle
  \qquad
  \text{for } \lambda \ge \lambda(\theta).
\]
Since
\[
  d'_\lambda = \left(0, 2\sin\left(\frac{\pi}{5}\right)\right),
\]
the derivative on the left branch is
\begin{align*}
  L'_\theta(\lambda)
  &=
  \langle -d'_\lambda, w_0(\theta)\rangle
  +
  \langle d'_\lambda, w_3(\theta)\rangle \\
  &=
  2\sin\left(\frac{\pi}{5}\right)
  \left(
    \sin\left(\theta + \frac{6\pi}{5}\right) - \sin\theta
  \right) \\
  &=
  4\sin\left(\frac{\pi}{5}\right)\sin\left(\frac{3\pi}{5}\right)
  \cos\left(\theta + \frac{3\pi}{5}\right)
  <
  0,
\end{align*}
because
$\theta + \frac{3\pi}{5} \in [\frac{3\pi}{5}, \frac{7\pi}{10}] \subset
(\frac{\pi}{2}, \pi)$.
On the right branch one gets
\begin{align*}
  L'_\theta(\lambda)
  &=
  \langle -d'_\lambda, w_0(\theta)\rangle
  +
  \langle d'_\lambda, w_2(\theta)\rangle \\
  &=
  2\sin\left(\frac{\pi}{5}\right)
  \left(
    \sin\left(\theta + \frac{4\pi}{5}\right) - \sin\theta
  \right) \\
  &=
  4\sin\left(\frac{\pi}{5}\right)\sin\left(\frac{2\pi}{5}\right)
  \cos\left(\theta + \frac{2\pi}{5}\right)
  \ge
  0,
\end{align*}
with strict inequality on $0 \le \theta < \pi/10$. Thus
$L_\theta(\lambda)$ decreases up to $\lambda(\theta)$ and increases
afterwards, so the unique minimizer is $\lambda(\theta)$.

\paragraph{Capacity along the active branch.}
Set
\[
  A := 1+\cos\left(\frac{\pi}{5}\right).
\]
The support-switch equation implies
\[
  (2\lambda(\theta)-1)\sin\left(\frac{\pi}{5}\right)
  =
  -A\tan\theta,
\]
so the minimizing chord vector is
\[
  d_{\lambda(\theta)} = (-A,-A\tan\theta).
\]
Therefore
\[
  h_{R(\theta)P_5}(-d_{\lambda(\theta)})
  =
  \langle -d_{\lambda(\theta)}, w_0(\theta)\rangle
  =
  A(\cos\theta+\tan\theta\sin\theta)
  =
  A\sec\theta.
\]
At the minimizing parameter one may use $w_2(\theta)$ for the forward segment,
because the support values at $w_2(\theta)$ and $w_3(\theta)$ agree there. Thus
\begin{align*}
  h_{R(\theta)P_5}(d_{\lambda(\theta)})
  &=
  \langle d_{\lambda(\theta)}, w_2(\theta)\rangle \\
  &=
  A\left(
    \cos\left(\frac{\pi}{5}\right)\cos\theta
    +
    \sin\left(\frac{\pi}{5}\right)\sin\theta
    +
    \tan\theta
    \left(
      \cos\left(\frac{\pi}{5}\right)\sin\theta
      -
      \sin\left(\frac{\pi}{5}\right)\cos\theta
    \right)
  \right) \\
  &=
  A\cos\left(\frac{\pi}{5}\right)\sec\theta.
\end{align*}
Hence the 2-bounce capacity candidate is
\[
  c_{2\text{-bounce}}(K_\theta)
  =
  A\left(1+\cos\left(\frac{\pi}{5}\right)\right)\sec\theta
  =
  A^2 \sec\theta
  \quad \text{for } 0 \le \theta \le \frac{\pi}{10}.
\]
At $\theta=0$, this candidate value is $A^2$.

The volume stays constant:
\[
  \operatorname{vol}(K_\theta)=\operatorname{area}(P_5)^2,
  \qquad
  \operatorname{area}(P_5)=\frac{5}{2}\sin\left(\frac{2\pi}{5}\right).
\]
Therefore
\[
  \mathrm{sys}(K_\theta)
  \le
  \frac{A^4}{2\operatorname{area}(P_5)^2}\sec^2\theta.
\]
Using
\[
  \sin^2\left(\frac{2\pi}{5}\right)=\frac{5+\sqrt{5}}{8},
  \qquad
  A=\frac{5+\sqrt{5}}{4},
\]
we obtain
\[
  \frac{A^4}{2\operatorname{area}(P_5)^2}
  =
  \frac{(5+\sqrt{5})^4}{256}
  \cdot
  \frac{16}{25(5+\sqrt{5})}
  =
  \frac{(5+\sqrt{5})^3}{400}
  =
  \frac{5+2\sqrt{5}}{10}.
\]
This proves the claimed secant law and the resulting upper bound on the
systolic ratio.
\end{proof}
\end{unverified}

\begin{unverified}
\begin{lemma}[Minimizing parameter for the active 2-bounce family]
\label{lem:pentagon-rotation-lambda}
For the candidate 2-bounce branch on $0 \le \theta \le \frac{\pi}{10}$, the
minimizing opposite-edge parameter is
\[
  \lambda(\theta)
  =
  \frac{1}{2}
  -
  \frac{\tan(\theta)}{2\tan(\pi/10)}.
\]
\end{lemma}
\end{unverified}

\begin{unverified}
\begin{lemma}[Capacity law on the active 2-bounce branch]
\label{lem:pentagon-rotation-capacity}
On the same branch,
\[
  c_{2\text{-bounce}}(K_\theta)
  =
  c_{2\text{-bounce}}(K_0)\sec(\theta).
\]
\end{lemma}
\end{unverified}

\begin{unverified}
\begin{lemma}[Empirical affine-branch picture]
\label{lem:pentagon-rotation-empirical-branch}
The owned sweep
`experiments/sys-landscape/pentagon-rotation-formula/theta-sweep.jsonl`
supports the following picture:
\begin{enumerate}
\item On the open half-domain $0 \le \theta < \pi/10$, every sampled minimum
  belongs to one affine 2-bounce class.
\item At $\theta=\pi/10$, many 2-bounce and 3-bounce minima tie.
\item On the second half-domain $\pi/10 < \theta \le \pi/5$, the minima fall
  into one mirrored affine 2-bounce class.
\end{enumerate}
This is a sanity check only. It is not part of the proof.
\end{lemma}
\end{unverified}

\begin{unverified}
\begin{lemma}[Strict-feasibility 3-bounce branches are smooth]
\label{lem:pentagon-rotation-three-bounce-smooth}
Fix a 3-bounce cyclic permutation~$\sigma$ from the billiard enumeration and
define
\[
  U_\sigma
  :=
  \left\{
    \theta \in [0,\pi/10] :
    M_\sigma(\theta) \text{ is non-singular},\;
    \beta_\sigma(\theta) > 0,\;
    Q_\sigma(\theta) > 0
  \right\}.
\]
Then $U_\sigma$ is open in~$[0,\pi/10]$. On each connected component
$I \subset U_\sigma$, the per-orbit action
\[
  A_\sigma(\theta) := \frac{1}{2Q_\sigma(\theta)}
\]
is smooth, hence continuous.
\end{lemma}

\begin{proof}
Let $a(\theta)$ denote the dual-vertex data of~$K_\theta$.
The map $\theta \mapsto a(\theta)$ is smooth because the second factor is
obtained from a fixed pentagon by rotation.
If $\theta_0 \in U_\sigma$, then
$M_\sigma(a(\theta_0))$ is non-singular,
$\beta_\sigma(a(\theta_0)) > 0$,
and $Q_\sigma(a(\theta_0)) > 0$.
Lemma~\ref{lem:orbit-feasibility-open} yields an $\varepsilon > 0$ such that
the same strict inequalities persist for
$|\theta-\theta_0|<\varepsilon$, so $U_\sigma$ is open.
Lemma~\ref{lem:per-orbit-smooth} applied to the smooth parameter
$a(\theta)$ shows that $A_\sigma(\theta)$ is smooth on the same
neighborhood. Since this holds at every point of~$U_\sigma$, the action is
smooth on each connected component~$I \subset U_\sigma$.
\end{proof}
\end{unverified}

\begin{unverified}
\begin{lemma}[Continuity reduction on a strict-feasibility interval]
\label{lem:pentagon-rotation-three-bounce-continuity}
Let~$\sigma$ be a 3-bounce cyclic permutation and let $I \subset U_\sigma$ be
a connected component as in
Lemma~\ref{lem:pentagon-rotation-three-bounce-smooth}. Define
\[
  g(\theta)
  :=
  \left(1+\cos\left(\frac{\pi}{5}\right)\right)^2 \sec(\theta),
\]
the explicit 2-bounce capacity from
Proposition~\ref{prop:pentagon-rotation-two-bounce}. If
\[
  A_\sigma(\theta) > g(\theta)
  \qquad
  \text{for almost every } \theta \in I,
\]
then
\[
  A_\sigma(\theta) \ge g(\theta)
  \qquad
  \text{for every } \theta \in I.
\]
\end{lemma}

\begin{proof}
By Lemma~\ref{lem:pentagon-rotation-three-bounce-smooth},
$A_\sigma$ is continuous on~$I$, and $g$ is continuous on~$[0,\pi/10]$.
If there were $\theta_0 \in I$ with $A_\sigma(\theta_0) < g(\theta_0)$, then
continuity of $A_\sigma-g$ would imply $A_\sigma(\theta) < g(\theta)$ on an
open neighborhood of~$\theta_0$ inside~$I$. That neighborhood has positive
measure, contradicting the assumed almost-everywhere strict inequality.
\end{proof}
\end{unverified}

\begin{unverified}
\begin{lemma}[Reduction of the 3-bounce blocker to branch interiors and endpoints]
\label{lem:pentagon-rotation-three-bounce-reduction}
Assume the following.
\begin{enumerate}
\item For every 3-bounce cyclic permutation~$\sigma$ and every connected
  component $I \subset U_\sigma$, one has
  $A_\sigma(\theta) > g(\theta)$ for almost every $\theta \in I$, where
  $g(\theta)$ is the explicit 2-bounce capacity from
  Lemma~\ref{lem:pentagon-rotation-three-bounce-continuity}.
\item For every admissible 3-bounce orbit at an angle
  $\theta_0 \in [0,\pi/10]$ that does not lie in any such strict-feasibility
  component, either:
  \begin{enumerate}
  \item some $\beta_k(\theta_0)=0$, so the orbit contracts to a 2-bounce orbit
    with the same action by Lemma~\ref{lem:orbit-contraction}, and that
    contracted 2-bounce orbit is shown to have action at least~$g(\theta_0)$;
    or
  \item that orbit is shown separately not to minimize.
  \end{enumerate}
\end{enumerate}
Then every minimizing 3-bounce trajectory on $0 \le \theta \le \pi/10$ has
$T^\circ$-length at least that of the minimizing 2-bounce family.
\end{lemma}

\begin{proof}
Fix an admissible 3-bounce orbit at angle~$\theta_0$ and choose a cyclic
permutation~$\sigma$ representing it.
If $\theta_0 \in U_\sigma$, then $\theta_0$ lies in a connected component
$I \subset U_\sigma$, and
Lemma~\ref{lem:pentagon-rotation-three-bounce-continuity} gives
$A_\sigma(\theta_0) \ge g(\theta_0)$.
If $\theta_0 \notin U_\sigma$, then assumption~(2) applies.
In case~(a), Lemma~\ref{lem:orbit-contraction} reduces the action to a
2-bounce orbit with the same value, and the remaining comparison with
$g(\theta_0)$ is part of assumption~(2a).
In case~(b), the orbit is not minimizing by assumption.
Thus no minimizing 3-bounce orbit can have action below~$g(\theta)$ on the
half-domain.
\end{proof}
\end{unverified}

\begin{unverified}
\begin{lemma}[3-bounce exclusion: current blocker]
\label{lem:pentagon-rotation-three-bounce}
Every minimizing 3-bounce trajectory on $0 \le \theta \le \pi/10$ has
$T^\circ$-length at least that of the minimizing 2-bounce family above.
\end{lemma}

% [GAP - Lemma~\ref{lem:pentagon-rotation-three-bounce-reduction} isolates the
% remaining burden. It is still open to verify the interior inequality
% $A_\sigma > g$ on each strict-feasibility interval for the relevant
% 3-bounce combinatorics, and to show that every admissible endpoint away from
% the midpoint is either a contraction point or cannot minimize. The natural
% candidate route is a theta-uniform replacement for the HKO midpoint sliding
% sign computation, but the active support labels move with theta and the
% midpoint equality analysis does not transplant verbatim.]
\end{unverified}

\begin{unverified}
\begin{remark}
The current computational evidence supports the theorem statement to
floating-point precision on the owned sweep
`experiments/sys-landscape/pentagon-rotation-formula/theta-sweep.jsonl`.
This is a sanity check only; it is not part of the proof.
\end{remark}
\end{unverified}
